
\section{Trust in the Congregations}

Editorial article in “Folkebladet” for April 26, 1899.

It is clearly noticeable throughout the whole ecclesiastical discussion among us that the urge to build “church bodies” derives, at least in large part, from a kind of distrust of the congregations. For what among us is called “church body” has hitherto, at least, become nothing other than organizations of pastors, in which the congregations as such have nothing to say. These organizations of pastors are regarded as trustworthy in the cause of religion; but people seem to be afraid to entrust the cause of Christianity, or at least of Lutheranism, to the congregations.

It is feared that the congregations either lack sufficient interest or sufficient insight, or both.

And there is surely some ground for this fear, since it is so common. Nor is it strange if people who come out of the coercion of the state church, and often without any living Christianity, should begin by supporting pastor and church, and, at least in the beginning, have both little interest and little understanding of the matter.

When one has never asked the people whether they want a congregation, or whether they want to be a congregation, but only whether they want a pastor and administrations of the sacraments, then it is indeed reasonable enough that many feel as though pastor and church and ceremonies and customs are a burden, which is necessary, but which they would prefer to see made as light as possible.

In other words, the distrust of the congregations, which is discernible among many pastors, perhaps is in reality a distrust of the work that has been done for the congregation’s edification among us.

If a shipbuilder does not dare to go aboard a ship which he himself has built, then that is not any especially strong recommendation of the shipbuilder himself.

If, however, we go to God’s Word, then we do not find this distrust of the congregations. The New Testament presents it in such a way that the congregation is equipped by God Himself with all the necessary gifts to preserve and spread God’s Kingdom on earth. Jesus and His apostles are not afraid—under God—to entrust the cause of the Gospel in the world to the congregation. There is no higher organization above the congregation that is more trustworthy than the congregation itself.

Is it to be understood, then, that the congregations in those days had greater interest and more insight? And did not the congregations’ competence grow by the responsibility that rested upon them?

When we think over this matter, it appears as though the congregations among us have lost their freedom because they have lost their competence. And now they are kept in tutelage, and thereby their incapacity continues.

Is it not then right for the Free Church to work for this matter to return once more to its proper course? And is it not fitting that freedom and responsibility, and responsibility and freedom, should go together, so that the one grows with the other?

Let the congregations receive the work and the responsibility that God’s Word gives them, and the cause of Christianity shall not suffer if it is entrusted to those to whom God’s Word entrusts it.

But do not misunderstand us. We do not mean a congregation without a pastor; when we speak of congregations, everything that belongs to the congregation is included therein. And we believe that it would be an exceedingly great advance if on both sides it could be acknowledged that there is no advantage in cultivating opposition between pastors and laypeople, since both are members of the congregation and need one another.

Georg Sverdrup
